% For printing in a4
\documentclass[a4,10pt,twoside,openright,italian,english]{book}% twoside!

% For printing with the A5 format
%\documentclass[10pt,twoside,openright,english,italian]{book}% twoside!

% Set paper size
\usepackage[twoside=true]{geometry}

%For printing with the weird format
%\geometry{
%	paperwidth=17cm,
%	paperheight=24cm,
%	margin=2cm,
%	top=2.3cm,
%	bindingoffset=0.4cm
%}
% For printing in a4
\geometry{a4paper,
  margin=3cm,
  top=3.8cm,
  bindingoffset=0.4cm
}

%Uncomment this for final prints: this just enables printing on a4 paper
\usepackage[cam,center,a4,pdflatex,axes]{crop}

\usepackage{phdthesis}

\usepackage{fancyhdr}
\usepackage{color}
\usepackage{array}
\usepackage{mdwmath}
\usepackage{mdwtab}
\usepackage{amsmath,amssymb}
\usepackage{cite}
\usepackage{graphicx}
\usepackage{listings}
\usepackage{subfig}
\usepackage{booktabs}
\usepackage{latexsym}
\usepackage{color}
\usepackage{url}
\usepackage{bnf}
\usepackage{rotating}
\usepackage{multirow}
\usepackage{phdtitle}
\usepackage{paralist}
\usepackage{bibentry}
%\usepackage[algochapter]{algorithm2e}
\usepackage[bookmarks=true,
%pdftex=false,
bookmarksopen=true,hidelinks]{hyperref}
\usepackage[toc,acronym]{glossaries}
\usepackage{lscape}
\usepackage{algorithmic}
\usepackage{algorithm}
\usepackage{longtable}
\usepackage[T1]{fontenc} 
%\usepackage[latin1]{inputenc}
\usepackage[utf8]{inputenc}
%\usepackage{fontspec}
%\setmainfont{Calibri}
%% \hyphenation{} is used to force the 
\hyphenation{}

\newtheorem{Definition}{Definition}[section]

\lstset{tabsize=2,basicstyle=\footnotesize,breaklines=true}

\usepackage[english,italian]{babel}
\usepackage{ulem}
\normalem

\hypersetup{pdftitle={Title}, pdfauthor={Author}}

\nobibliography*

\department {Dottorato di ricerca in Ingegneria dell'Informazione}

% Please fulfil the followed fields with your data
\author{Name Surname} 
\title{Title}
\tutor{Tutor Name}
\supervisor{Coordinator Name}
\titleimage{img/unipi.png} %please not change
\phdyear{Year}
\phdmonth{Month}
\phdcycle{Cycle}

%%%%%%%%%%%%%%%%%%%%%%%%%%%%%%%%%%%%%%
% Let's Start The Real Document
%%%%%%%%%%%%%%%%%%%%%%%%%%%%%%%%%%%%%%
% \newglossaryentry{matrix_channel}
% {       name={$H^*$},
%         description={Conjugate operation}
% }
% \newglossaryentry{trasp_x}
% {       name={$[ x ]^{\rm T}$},
%         description={transpose operator}
% }
% \newglossaryentry{vec_x}
% {       name={\textbf{x}},
%         description={vectors are in bold}
% }
% \newglossaryentry{floor_funct}
% {       name={$\left\lfloor x \right\rfloor$},
%         description={round to the lower integer of $x$}
% }

% PREFIXES

\newglossaryentry{ogc}
{       name={ogc:},
        description={\texttt{http://www.opengis.net/}}
}

\newglossaryentry{geo}
{       name={geo:},
        description={\texttt{http://www.opengis.net/ont/geosparql#}}
}

\newglossaryentry{geof}
{       name={geof:},
        description={\texttt{http://www.opengis.net/def/function/geosparql/}}
}

\newglossaryentry{geor}
{       name={geor:},
        description={\texttt{http://www.opengis.net/def/rule/geosparql/}}
}

\newglossaryentry{rdf}
{       name={rdf:},
        description={\texttt{http://www.w3.org/1999/02/22-rdf-syntax-ns#}}
}

\newglossaryentry{rdfs}
{       name={rdfs:},
        description={\texttt{http://www.w3.org/2000/01/rdf-schema#}}
}

\newglossaryentry{xsd}
{       name={xsd:},
        description={\texttt{http://www.w3.org/2001/XMLSchema#}}
}

\newglossaryentry{owl}
{       name={owl:},
        description={\texttt{http://www.w3.org/2002/07/owl#}}
}

\newglossaryentry{schema}
{       name={schema:},
        description={\texttt{http://schema.org/}}
}

\newglossaryentry{strdf}
{       name={strdf:},
        description={\texttt{http://strdf.di.uoa.gr/ontology#}}
}

\newglossaryentry{ecrm}
{       name={ecrm:},
        description={\texttt{http://erlangen-crm.org/current/}}
}

\newglossaryentry{crmgeo}
{       name={crmgeo:},
        description={\texttt{http://www.ics.forth.gr/isl/CRMgeo/}}
}

\newglossaryentry{narra}
{       name={narra:},
        description={\texttt{https://dlnarratives.eu/ontology#}}
}

\newglossaryentry{igen}
{       name={igen:},
        description={\texttt{https://imagoarchive.it/thes/tid/}}
}

\newglossaryentry{ilrm}
{       name={ilrm:},
        description={\texttt{http://imagoarchive.it/ilrmoo/}}
}

\newglossaryentry{wdt}
{       name={wdt:},
        description={\texttt{http://www.wikidata.org/prop/direct/}}
}
\newglossaryentry{wd}
{       name={wd:},
        description={\texttt{http://www.wikidata.org/entity/}}
}

\newglossaryentry{uom}
{       name={uom:},
        description={\texttt{http://www.opengis.net/def/uom/OGC/1.0/}}
}

\newglossaryentry{osm}
{       name={osm:},
        description={\texttt{https://www.openstreetmap.org/}}
}

\newglossaryentry{osm2rdfkey}
{       name={osm2rdfkey:},
        description={\texttt{https://osm2rdf.cs.uni-freiburg.de/rdf/key#}}
}

\newglossaryentry{osmkey}
{       name={osmkey:},
        description={\texttt{https://www.openstreetmap.org/wiki/Key:}}
}

\newglossaryentry{imago}
{       name={imago:},
        description={\texttt{https://imagoarchive.it/ontology/}}
}










% ogc: http://www.opengis.net/
% geo: http://www.opengis.net/ont/geosparql#
% geof: http://www.opengis.net/def/function/geosparql/
% geor: http://www.opengis.net/def/rule/geosparql/
% sf: http://www.opengis.net/ont/sf#
% gml: http://www.opengis.net/ont/gml#
% xsd: http://www.w3.org/2001/XMLSchema#
% rdf: http://www.w3.org/1999/02/22-rdf-syntax-ns#
% rdfs: http://www.w3.org/2000/01/rdf-schema#
% owl: http://www.w3.org/2002/07/owl#
\newacronym[longplural={geographic information systems}]{GISLabel}{GIS}{geographic information systems}
\newacronym[longplural={Resource Description Framework}]{RDFLabel}{RDF}{Resource Description Framework}
\newacronym[longplural={Resource Description Framework Schema}]{RDFSLabel}{RDFS}{Resource Description Framework Schema}
\newacronym[longplural={Web Ontology Language}]{OWLLabel}{OWL}{Web Ontology Language}
\newacronym[longplural={SPARQL Protocol and RDF Query Language }]{SPARQLLabel}{SPARQL}{SPARQL Protocol and RDF Query Language }
\newacronym[longplural={Uniform Resource Identifiers}]{URILabel}{URI}{Uniform Resource Identifier}
\newacronym[longplural={Internationalized Resource Identifiers}]{IRILabel}{IRI}{Internationalized Resource Identifier}
\newacronym[longplural={Linked Open Data}]{LODLabel}{LOD}{Linked Open Data}
\newacronym[longplural={Terse RDF Triple Language}]{TurtleLabel}{Turtle}{Terse RDF Triple Language}
\newacronym[longplural={Coordinate Reference Systems}]{CRSLabel}{CRS}{Coordinate Reference System}
\newacronym[longplural={Region Connection Calculus}]{RCCLabel}{RCC}{Region Connection Calculus}
\newacronym[longplural={jointly exhaustive and pairwise disjoint}]{JEPDLabel}{JEPD}{jointly exhaustive and pairwise disjoint}
\newacronym[longplural={9-intersection model}]{9IMLabel}{9IM}{9-intersection model}
\newacronym[longplural={dimensionally extended 9-intersection model}]{DE9IMLabel}{DE9IM}{dimensionally extended 9-intersection model}
\newacronym[longplural={4-intersection model}]{4IMLabel}{4IM}{4-intersection model}
\newacronym[longplural={Open Geospatial Consortium}]{OGCLabel}{OGC}{Open Geospatial Consortium}
\newacronym[longplural={Event Calculus}]{ECLabel}{EC}{Event Calculus}
\newacronym[longplural={Descriptive Ontology for Linguistic and Cognitive Engineering}]{DOLCELabel}{DOLCE}{Descriptive Ontology for Linguistic and Cognitive Engineering}
\newacronym[longplural={Well-Known Text}]{WKTLabel}{WKT}{Well-Known Text}
\newacronym[longplural={Geography Markup Language}]{GMLLabel}{GML}{Geography Markup Language}
\newacronym[longplural={International Committee for Documentation}]{CIDOCLabel}{CIDOC}{International Committee for Documentation}
\newacronym[longplural={Conceptual Reference Model}]{CRMLabel}{CRM}{Conceptual Reference Model}
\newacronym[longplural={CIDOC Conceptual Reference Model}]{CIDOCCRMLabel}{CIDOC CRM}{CIDOC Conceptual Reference Model}
\newacronym[longplural={Narrative Ontology}]{NOntLabel}{NOnt}{Narrative Ontology}
\newacronym[longplural={Geospatial Narrative Ontology}]{NOnt+SLabel}{NOnt+S}{Geospatial Narrative Ontology}
\newacronym[longplural={Semantic Web Rule Language}]{SWRLLabel}{SWRL}{Semantic Web Rule Language}
\newacronym[longplural={Value Chains}]{VCLabel}{VC}{Value Chain}
\newacronym[longplural={Local Administrative Units}]{LAULabel}{LAU}{Local Administrative Unit}
\newacronym[longplural={Nomenclature of Territorial Units for Statistics}]{NUTSLabel}{NUTS}{Nomenclature of Territorial Units for Statistics}
\newacronym[longplural={Findable, Accessible, Interoperable, and Reusable}]{FAIRLabel}{FAIR}{Findable, Accessible, Interoperable, and Reusable}
\newacronym[longplural={Story Map Building And Visualising Tool}]{SMBVTLabel}{SMBVT}{Story Map Building And Visualising Tool}
\newacronym[longplural={LinkedGeoData}]{LGDLabel}{LGD}{LinkedGeoData}
\newacronym[longplural={OpenStreetMap}]{OSMLabel}{OSM}{OpenStreetMap}
\newacronym[longplural={Greek Administrative Geography}]{GAGLabel}{GAG}{Greek Administrative Geography}
\newacronym[longplural={CORINE Land Cover}]{CLCLabel}{CLC}{CORINE Land Cover}
\newacronym[longplural={Earth Observation}]{EOLabel}{EO}{Earth Observation}
\newacronym[longplural={Non-Topological Construct Functions}]{NTCFLabel}{NTCF}{Non-Topological Construct Functions}
\newacronym[longplural={Spatial Selections}]{SSLabel}{SS}{Spatial Selections}
\newacronym[longplural={Spatial Joins}]{SJLabel}{SJ}{Spatial Joins}
\newacronym[longplural={Geographic Information System of the Commission}]{GISCOLabel}{GISCO}{Geographic Information System of the Commission}
\newacronym[longplural={MOuntain Valorisation through INterconnectedness and Green growth}]{MOVINGLabel}{MOVING}{MOuntain Valorisation through INterconnectedness and Green growth}
\newacronym[longplural={Index Medii Aevi Geographiae Operum}]{IMAGOLabel}{IMAGO}{Index Medii Aevi Geographiae Operum}


\makeglossaries

\begin{document}
\selectlanguage{english}

\maketitle

\pagestyle{empty}

\cleardoublepage
\newpage

%%%%%%%%%%%%%%%%%%%%%%%%%%%%%%%%%%%%%%
% Dedication - For removing dedication, please comment or delete the code inside \thispagestyle{empty}
%%%%%%%%%%%%%%%%%%%%%%%%%%%%%%%%%%%%%%
\thispagestyle{empty}
    \null\vspace{\stretch {1}}
        \begin{flushright}
                This thesis is dedicated to....
        \end{flushright}
\vspace{\stretch{2}}\null

\cleardoublepage
\newpage

\pagestyle{empty}
%% change numbering into Roman numbers for the introductory part
%\setcounter{page}{1}
%\pagenumbering{Roman}

%%%%%%%%%%%%%%%%%%%%%%%%%%%%%%%%%%%%%%
% Quotes - For removing quotes, please comment or delete the code inside \thispagestyle{empty}
%%%%%%%%%%%%%%%%%%%%%%%%%%%%%%%%%%%%%%
\thispagestyle{empty}
    \null\vspace{\stretch {1}}
        \begin{flushright}
                "A few first rate research papers are preferable to a large number \\
                that are poorly conceived or half-finished.\\
                The latter are no credit to their writers and \\
                a waste of time to their readers"\\
                Claude Shannon
                %IRE Transactions on Information Theory (1956), volume 2, issue 1, page 3. Shannon, Claude E. (March 1956), The Bandwagon, 2, doi:10.1109/TIT.1956.1056774.
        \end{flushright}
\vspace{\stretch{2}}\null

\cleardoublepage
\newpage

\pagestyle{empty}
%% change numbering into Roman numbers for the introductory part
\setcounter{page}{1}
\pagenumbering{Roman}

%%%%%%%%%%%%%%%%%%%%%%%%%%%%%%%%%%%%%%
% Acknowledgement
%%%%%%%%%%%%%%%%%%%%%%%%%%%%%%%%%%%%%%
\chapter*{Acknowledgements}
\lettrine{A}{cknowledgements} goes here.

\selectlanguage{italian}
\chapter*{Ringraziamenti}
\lettrine{R}{ingraziamenti} 
\selectlanguage{english}

\cleardoublepage
\newpage

\pagestyle{fancy}
% change numbering into Roman numbers for the introductory part

%%%%%%%%%%%%%%%%%%%%%%%%%%%%%%%%%%%%%%
% Summary
%%%%%%%%%%%%%%%%%%%%%%%%%%%%%%%%%%%%%%
\selectlanguage{english}
\chapter*{Summary}

\lettrine{G}{eospatial} knowledge plays a fundamental role in shaping the understanding and communication of narratives. Whether historical, literary, or journalistic, narratives often incorporate spatial elements that provide essential context, enriching the reader’s engagement with the story. This thesis addresses a fundamental challenge in the fields of knowledge representation and the Semantic Web: the effective semantic modeling of geospatial narratives. Geospatial narratives are conceived as networks of events interconnected through semantic links, capturing the complex relationships between spatial occurrences. Current models lack sufficient expressiveness and interoperability, which limits their capacity for data integration, querying, and analysis. To overcome these limitations, this research develops the Geospatial Narrative Ontology (NOnt+S), an extension of the existing Narrative Ontology (NOnt) that is based on the ISO standard CIDOC Conceptual Reference Model (CIDOC CRM). By formally modeling geospatial information using a formal language like OWL 2 DL, the ontology ensures decidability and enhances interoperability across diverse knowledge bases. Furthermore, implementing this ontology within a Semantic Web framework and employing semantic reasoners enable the inference of new geospatial knowledge, thereby enriching data analysis and facilitating more sophisticated queries. Finally, the thesis addresses the challenges of visualization by proposing Story Maps as a solution to visualize geospatial narratives. Story Maps are interactive digital maps that tell stories by integrating geospatial data with narrative elements. Story mapping, in this context, is a method for arranging narrative events to create a more holistic view of how they fit into the overall geospatial narrative, thereby enhancing user engagement and understanding.

\section*{Research Objectives}

The thesis aims to advance geospatial knowledge representation in narratives by:
\begin{enumerate}
    \item Developing an ontology, NOnt+S, wich representan geospatial and temporal aspects of narratives ensuring interoperability with existing Semantic Web frameworks.
    \item Enhancing data integration and enabling advanced geospatial querying mechanisms to achieve consistent and comprehensive representations of geospatial information in narratives.
    \item Evaluating the ontology's practical implications and applicability, including its ability to infer new knowledge, through empirical validation in real-world case studies from the MOVING (Horizon 2020-2024) and IMAGO (PRIN 2020-2024) projects.
    \item Developing effective visualization methods for geospatial narratives by employing Story Maps, integrating geospatial data with narrative elements to enhance user understanding and engagement.
\end{enumerate}

\section*{Methodology}

The research employs a rigorous, structured methodology based on the METHONTOLOGY framework, which includes:
\begin{itemize}
    \item \textbf{Requirements Specification}: Identifying the requirements for geospatial narrative representation through comprehensive analysis.
    \item \textbf{Conceptualization and Logical Formalization}: Extending the NOnt ontology to incorporate geospatial and temporal constructs, providing formal mathematical specifications.
    \item \textbf{Implementation}: Developing NOnt+S in OWL 2 DL, ensuring decidability and interoperability to achieve syntactic precision and compatibility within the Semantic Web environment.
    \item \textbf{Evaluation}: Validating the ontology through case studies, assessing its performance in terms of semantic enrichment, reasoning capabilities, and querying effectiveness.
\end{itemize}

\section*{Significance of Findings}
This thesis advances Knowledge Representation by introducing the NOnt+S ontology, a robust framework for modeling complex geospatial narrative structures. NOnt+S integrates narrative elements like fabula (chronological sequence) and plot (narrative presentation) with spatial constructs such as qualitative and quantitavce place. By bridging narrative and spatial dimensions, it addresses limitations in existing ontologies and sets a new standard for flexible modeling applicable in digital humanities, GIS, and smart city applications.

Adopting established standards like the CIDOC CRM and GeoSPARQL ensures interoperability with existing data systems and semantic web technologies. An embedded semantic reasoner allows NOnt+S to infer new knowledge from existing data, uncovering hidden patterns within narrative and spatial information.

The study also developed a software architecture utilizing Story Maps, providing an intuitive interface for visualizing geospatial narratives by integrating text, images, and maps. By supporting the construction of structured geospatial knowledge graphs and aligning with Linked Open Data principles, NOnt+S promotes data reusability and interdisciplinary collaboration. Validated across diverse domains, it demonstrates effectiveness in managing intricate geospatial queries and serves as a significant asset for advanced geospatial narrative modeling, laying the groundwork for future advancements.


\selectlanguage{english}

%%%%%%%%%%%%%%%%%%%%%%%%%%%%%%%%%%%%%%
% Italian Summary
%%%%%%%%%%%%%%%%%%%%%%%%%%%%%%%%%%%%%%

%There must be an Italian version of the summary
\selectlanguage{italian}
\chapter*{Sommario}
\lettrine{L}{a} conoscenza geospaziale ha un impatto significativo sulla comprensione e sulla comunicazione delle storie. Le narrazioni — che siano di natura storica, letteraria o giornalistica — spesso includono elementi spaziali che offrono un contesto fondamentale per la comprensione del racconto. Questa tesi affronta una sfida fondamentale nei campi della rappresentazione della conoscenza e del Web Semantico: la modellazione semantica delle narrazioni geospaziali. Le narrazioni geospaziali sono concepite come reti di eventi interconnessi attraverso legami semantici, catturando le complesse relazioni tra le occorrenze spaziali. I modelli attuali mancano di sufficiente espressività e interoperabilità, il che limita la loro capacità di integrazione dei dati, interrogazione e analisi. Per superare queste limitazioni, questa ricerca sviluppa l'Ontologia delle Narrazioni Geospaziali (NOnt+S), un'estensione dell'esistente Ontologia delle Narrazioni (NOnt) che si basa sullo standard ISO CIDOC Conceptual Reference Model (CIDOC CRM). La modellazione formale delle informazioni geospaziali tramite un linguaggio formale come OWL 2 DL garantisce la decidibilità e migliora l'interoperabilità tra diverse basi di conoscenza. Inoltre, implementando questa ontologia all'interno di un framework del Web Semantico e impiegando ragionatori semantici, si rende possibile l'inferenza di nuova conoscenza geospaziale, arricchendo così l'analisi dei dati e facilitando interrogazioni più sofisticate. Infine, la tesi affronta le sfida di una visualizzazione efficace proponendo le \textit{Story Maps} come soluzione per visualizzare le narrazioni geospaziali. Le \textit{Story Maps} sono mappe digitali interattive che raccontano storie integrando dati geospaziali con elementi narrativi. Lo \textit{story mapping}, in questo contesto, è un metodo per organizzare gli eventi narrativi al fine di creare una visione più olistica di come essi si inseriscono nella narrazione geospaziale complessiva, migliorando così il coinvolgimento e la comprensione dell'utente.

\section*{Obiettivi della ricerca}

La tesi mira ad avanzare la rappresentazione della conoscenza geospaziale nelle narrazioni proponendosi di:

\begin{enumerate}
    \item Sviluppare un'ontologia, NOnt+S, che rappresenta gli aspetti geospaziali e temporali delle narrazioni, garantendo l'interoperabilità con i framework esistenti del Web Semantico.
    \item Migliorare l'integrazione dei dati e abilitare sitemi di interrogazione geospaziale per ottenere rappresentazioni coerenti e complete delle informazioni geospaziali nelle narrazioni
    \item Valutare le implicazioni pratiche e l'applicabilità dell'ontologia, inclusa la sua capacità di inferire nuova conoscenza, attraverso la validazione empirica in casi di studio reali dai progetti MOVING (Horizon 2020-2024) e IMAGO (PRIN 2020-2024).
    \item Sviluppare metodi efficaci di visualizzazione per le narrazioni geospaziali utilizzando Story Maps, integrando dati geospaziali con elementi narrativi per migliorare la comprensione e il coinvolgimento dell'utente.
\end{enumerate}

\section*{Metodologia}

La ricerca impiega una metodologia rigorosa e strutturata basata sul framework METHONTOLOGY, che include:

\begin{itemize}
    \item \textbf{Specifica dei requisiti}: identificazione dei requisiti per la rappresentazione delle narrazioni geospaziali attraverso un'analisi esaustiva.
    \item \textbf{Concettualizzazione e formalizzazione logica}: estendere l'ontologia NOnt per incorporare costrutti geospaziali e temporali, fornendo una formalizzazione matematica.
    \item \textbf{Implementazione}: sviluppare NOnt+S in OWL 2 DL, garantendo decidibilità e interoperabilità per ottenere precisione sintattica e interoperabilità all'interno del Web Semantico.
    \item \textbf{Valutazione}: validare l'ontologia attraverso casi di studio, valutando le sue prestazioni in termini di arricchimento semantico, capacità di ragionamento ed efficacia nelle interrogazioni.
\end{itemize}

\section*{Importanza dei risultati}

Questa tesi avanza la rappresentazione della conoscenza introducendo l'ontologia NOnt+S, un framework robusto per modellare narrazioni geospaziali complesse. NOnt+S integra elementi narrativi come la fabula (sequenza cronologica) e l'intreccio (presentazione narrativa) con costrutti spaziali quali i luoghi qualitativi e quantitativi. Collegando le dimensioni narrative e spaziali, affronta le limitazioni delle ontologie esistenti e stabilisce un nuovo standard per la modellazione flessibile applicabile nelle \textit{digital humanities}, nei GIS e nelle applicazioni per le smart city.

Adottando standard consolidati come il CIDOC CCnceptual Reference Model (CIDOC CRM) e GeoSPARQL, si garantisce l'interoperabilità con i dati esistenti e le tecnologie del Web Semantico. Tramite un ragionatore semantico integrato, NOnt+S è in grado di inferire nuova conoscenza dai dati esistenti, rivelando pattern nascosti all'interno delle informazioni narrative e geospaziali.

Lo studio ha inoltre sviluppato un \textit{framework} software che utilizza Story Maps, fornendo un'interfaccia intuitiva per visualizzare narrazioni geospaziali attraverso l'integrazione di testo, immagini e mappe. Supportando la costruzione di grafi di conoscenza geospaziale strutturati e allineandosi ai principi dei Linked Open Data (LOD), NOnt+S promuove la riutilizzabilità dei dati e la collaborazione interdisciplinare. Validata in domini diversi, dimostra efficacia nella gestione di interrogazioni geospaziali complesse e rappresenta una risorsa significativa per la modellazione avanzata di narrazioni geospaziali, gettando le basi per futuri progressi.
\selectlanguage{english}

%%%%%%%%%%%%%%%%%%%%%%%%%%%%%%%%%%%%%%
% Publications
%%%%%%%%%%%%%%%%%%%%%%%%%%%%%%%%%%%%%%

%List of publications of the PhD candidate
\selectlanguage{english}
\chapter*{List of publications}

\section*{International Journals}
\begin{enumerate}
    % for citing your paper, take directly APA format from Google scholar
    % \item Surname, N., Surname, N. and Surname, N. (Year,Month). Title of the paper or journal. \emph{Place of publication}. (Vol. 123, pp. 456). EditorName.
    \item Bartalesi, V., Pratelli, N. (2024) Representing Geospatial Knowledge in Narratives.  \emph{ACM Journal on Computing and Cultural Heritage}, Just Accepted (November 2024).
    \item Bartalesi, V., Coro, G., Lenzi, E., Pratelli, N., Pagano, P., Moretti, M., Brunori, G. (2024). A Semantic Knowledge Graph of European Mountain Value Chains. \emph{Scientific Data}, 11(1), 978.
    \item Zabulis, X., Partarakis, N., Bartalesi, V., Pratelli, N., Meghini, C., Dubois, A., Moreno, I., Manitsaris, S. (2024). Multimodal Dictionaries for Traditional Craft Education. \emph{Multimodal Technologies and Interaction}, 8(7), 63.
    \item Zabulis, X., Partarakis, N., Demeridou, I., Bartalesi, V., Pratelli, N., Meghini, C., Nikolaou, N.,  Fallahian, P. (2024). Modelling and Simulation of Traditional Craft Actions. \emph{Applied Sciences}, 14(17), 7750.
    \item Bartalesi, V., Coro, G., Lenzi, E., Pagano, P., Pratelli, N. (2023). From unstructured texts to semantic story maps. \emph{International Journal of Digital Earth}, 16(1), 234–250.
    \item Bartalesi, V., Coro, G., Lenzi, E., Pratelli, N., Pagano, P., Felici, F., Moretti, M., Brunori, G. (2023). Using semantic story maps to describe a territory beyond its map. \emph{Semantic Web}, 14(6), 1255–1272.
    \item Bartalesi, V., Pratelli, N., Lenzi, E., Pontari, P. (2023). Using Semantic Web to Create and Explore an Index of Toponyms Cited in Medieval Geographical Works. \emph{ACM Journal on Computing and Cultural Heritage}, 16(2), 1–18.
    \item Zabulis, X., Partarakis, N., Demeridou, I., Doulgeraki, P., Zidianakis, E., Argyros, A., Theodoridou, M., Marketakis, Y., Meghini, C., Bartalesi, V. (2023). A roadmap for craft understanding, education, training, and preservation. \emph{Heritage}, 6(7), 5305–5328.
    \item Bartalesi, V., Metilli, D., Pratelli, N.,  Pontari, P. (2022). Towards a knowledge base of medieval and renaissance geographical Latin works: The IMAGO ontology. \emph{Digital Scholarship in the Humanities}, 37(1), 34–50.
    \item Bartalesi, V., Pratelli, N., Lenzi, E. (2022). Linking different scientific digital libraries in Digital Humanities: The IMAGO case study. \emph{International Journal on Digital Libraries}, 23(4), 303–317.
    \item Bartalesi, V., Pratelli, N., Meghini, C., Metilli, D., Tomazzoli, G., Livraghi, L. M., Zaccarello, M. (2022). A formal representation of the divine comedy’s primary sources: The Hypermedia Dante Network ontology. \emph{Digital Scholarship in the Humanities}, 37(3), 630–643.
\end{enumerate}

\section*{International Conferences/Workshops with Peer Review}
\begin{enumerate}
    % for citing your paper, take directly APA format from Google scholar
    % \item Surname, N., Surname, N. and Surname, N. (Year,Month). Title of the paper or journal. \emph{Place of publication}. (Vol. 123, pp. 456). EditorName.
    
    \item Bartalesi, V., Lenzi, E., Pratelli, N. (2023). A Web Tool to Create and Visualise Semantic Story Maps. 163–169.
    \item Bartalesi, V., Pratelli, N. (2021). The imago project: Towards a knowledge base of medieval and renaissance geographical works. SWODCH.
    \item Bartalesi, V., Pratelli, N., Lenzi, E. (2022). A Knowledge Base of Medieval and Renaissance Geographic Latin Works. IRCDL.
    \item Pratelli, N. (2022). A Geographical Extension for NOnt Ontology. 3246, 126–133.
\end{enumerate}

\section*{Others}
\begin{enumerate}
    % for citing your paper, take directly APA format from Google scholar
    % \item Surname, N., Surname, N. and Surname, N. (Year,Month). Title of the paper or journal. \emph{Place of publication}. (Vol. 123, pp. 456). EditorName.
    \item Bartalesi, V., Coro, G., Lenzi, E., Pratelli, N., Pagano, P. (2023). Towards Digital Twins of Territories Through Semantic Story Maps. \emph{BUILding a DIgital Twin: requirements, methods, and applications}.
    \item Bartalesi, V., Lenzi, E., Pratelli, N. (2023). Enhancing Narratives in Digital Humanities: Harnessing the Capabilities of Semantic Story Maps. Life Narrative and the Digital 2023.
\end{enumerate}















\selectlanguage{english}

%%%%%%%%%%%%%%%%%%%%%%%%%%%%%%%%%%%%%%
% List of Abbreviation and Symbols
%%%%%%%%%%%%%%%%%%%%%%%%%%%%%%%%%%%%%%
%Have a look at \gls command... you can refer a term in multiple ways
\printglossary[type=\acronymtype,title=List of Abbreviations]
\let\cleardoublepage\clearpage
\printglossary[title=Notation]
\cleardoublepage
\newpage
%%%%%%%%%%%%%%%%%%%%%%%%%%%%%%%%%%%%%%
% TOC
%%%%%%%%%%%%%%%%%%%%%%%%%%%%%%%%%%%%%%
\tableofcontents
\cleardoublepage
\newpage

% Now lets go back to normal numbering
\setcounter{page}{1}
\pagenumbering{arabic}

\cleardoublepage
\chapter{Title chapter 1}

%============================= INTRODUCTION =================================
\section{Introduction}\label{sec:i}
Remember that wherever you want you can cite with \verb.\cite{a1}. \cite{a1}, \cite{b1} and/or \cite{m1} someone.
Or use \verb.\footnote{That's a footnote}. like this\footnote{That's a footnote} 
% this is used to include other .tex files in order to split a long document 

%============================= SECTIONS =================================

\section{New section}
You can insert a definition
	\begin{Definition}
		Unipi: Universit\'a di Pisa
	\end{Definition}
	
\section{Another section}
	\subsection{And a subsection}
		In this sub section I will include many images just to make the list of figures meaningful.

			\begin{figure}[h!tb]
				\centerline {\includegraphics[scale=0.6]{img/unipi.png}}
				\caption{Caption of this Unipi image.}
				\label{fig:leet}
			\end{figure}

\section{New section}
	\subsection{And also many sub}
		\subsubsection{Or again more nested}
			\begin{itemize}
				\item item 1
				\item item 2
				\item item 3
			\end{itemize}			
			
\section{New section}
This is an example of table. You can see it in Table~\ref{tab:TableName}. 

 \begin{table}[h!tb]
   \centering \caption{Table caption}
   \label{tab:TableName}
   \vskip 0.2cm
   %%
   \scalebox{0.90}{
	    %% The {|c|c|c|c|c|} define the number of columns.
	    %% c means centered
	    %% | defines a vertical line between two columns 
	    \begin{tabular}{|c|c|c|}
	      \hline
	      Col 1 & Col 2 & Col2  \\
	      %% \\ force a newline without creating an horizontal line
	      Dim C.1 & Dim C.3 & Dim C.3 \\
	      %% \hline create an horizontal line between two rows
	      \hline
	      Data 1.1 & Data 1.2 & Data 1.3 \\
	      \hline
	      Data 2.1 & Data 2.2 & Data 2.3 \\
	      \hline
	      Data 3.1 & Data 3.2 & Data 3.3 \\  
	      \hline
	    \end{tabular}
	 }
 \end{table}

\section{You can do as many section as you want}
The following equation has been created using \verb.$$.: $ y_1 = a*x + b$.

%% \smallskip introduces a small vertical skip
\smallskip

%% \noindent is used to eliminate the blankline and the paragraph settings 
\noindent
This is a more complex equation: \[ y_2 = a*x + b \], created using \verb.\[ \].

%% \bigskip introduces a big vertical skip
\bigskip

\noindent
This is, (\ref{eq:eeq1}), an enumerated equation: 

	\begin{equation}\label{eq:eeq1}
		y_3 = a*x + b
	\end{equation}

, created using:

	\begin{verbatim}
		\begin{equation}\label{eq:eeq1}
		    y_3 = a*x + b
		\end{equation}
	\end{verbatim}

It is also possible to insert acronyms, after definition in \textit{Acronyms} file, like \acrshort{fpsLabel} which stands for \acrlong{fpsLabel} or \acrshort{SNRLabel} which stands for \acrlong{SNRLabel}.
Please define the notation at the beginning, defining operators and symbols in \textit{Notation} file, such as \gls{matrix_channel}, \gls{trasp_x}, \gls{vec_x}, \gls{floor_funct}. % uncomment if you want part I
% Insert further chapters using the following syntax
% \input{Source_Folder_Name/FileName}


\cleardoublepage
\phantomsection
\addcontentsline{toc}{chapter}{\bibname}
\small
\bibliographystyle{plain}
\bibliography{thesis}

\end{document}
